\begin{definition}
    For positive integers, $n \geq k$ and $r > s$, let $A_k(n; r, s)$ denote the least $N$ such that every $r$-coloring of $[N]^{(k)}$ contains a tight monotone path on $n$ edges which uses at most $s$ colors.
\end{definition}

A trivial bound for $A_{2}(n;3,2)$ comes from the observation that a $3$-coloring of $[N]^{(2)}$ can be regarded as a $2$-coloring of $[N]^{(2)}$ if one pretends that two of the colors are the same, and that a monochromatic path in the latter setting is a path which avoids the usage of a color in the original $3$-coloring. The Erd\H{o}s--Szekeres theorem then implies that
$$A_2(n; 3, 2) \leq MS_{2}(n;2) = n^2+1.$$
In \cite{Loh15}, Loh noted that $A_{2}(n;3,2) \geq n^{3/2}$ and used the proof of the Erd\H{o}s--Szekeres theorem due to Seidenberg \cite{Seid59} together with the triangle removal lemma of Ruzsa and Szemer\'edi \cite{RS78} to show that $A_{2}(n;3,2)=o(n^2)$ holds. In a remarkable paper \cite{GL21}, Gowers and Long further showed that there exists an absolute constant $\e > 0$ such that $A_2(n; 3, 2) \leq n^{2-\e}$ and conjectured that $A_{2}(n;3,2) = \Theta(n^{3/2})$. 


\subsection{Our results}

A more general first estimate on $A_{k}(n;r,s)$ can be obtained in a similar fashion from Theorem \ref{MS} by noticing that an $r$-coloring of $[N]^{(k)}$ can be regarded as a $\lceil r/s \rceil$-coloring of $[N]^{(k)}$ if one pretends that $\lceil r/s \rceil$ disjoint groups of $s$ colors each represent a single color, and that a monochromatic tight path in the latter setting is a tight path which uses at most $s$ colors in the original $r$-coloring. For example, when $k=3$, $r=3$, and $s=2$, this gives
$$A_3(n; 3, 2) \leq MS_3(n; 2) = \binom{2n}{n} + 1.$$
In general, for every $n \geq k \geq 3$, $r > s \geq 1$,
$$A_k(n; r, s) \leq MS_k(n; \lceil r/s \rceil) \leq T_{k-2}(2n^{\lceil r/s \rceil}).$$

Our main result is an improvement on the height of the tower in these upper bounds by one exponential, in the case when $k \geq 3$ and $s > r/2$. For example, when $k=3$, we show the following:

\begin{theorem}\label{main3}
    There exists an absolute constant $C > 0$ such that for all positive integers $n \geq 3$ and $r$, $s$ with $s > r/2$, 
    \begin{equation*}
        A_3(n; r, s) \leq n^{C \binom{r}{s}^2 \log \binom{r}{s}}.
    \end{equation*}
\end{theorem}

Moreover, for the particular combination $r=3$ and $s=2$, our proof gives the stronger bound $A_{3}(n;3,2) \leq n^{9}$. Qualitatively speaking, Theorem \ref{main3} shows that both $A_2(n;r,s)$ and $A_3(n;r,s)$ are polynomial in $n$ for $s > r/2$. We also show that for uniformity $k \geq 4$, the tower height continues to require one less exponential than the trivial bound from the theorem of Moshkovitz and Shapira.

\begin{theorem}\label{maink}
    There exists an absolute constant $C > 0$ such that for all positive integers $n \geq k \geq 3$ and $r$, $s$ with $s > r/2$, we have
    \begin{equation*}
        A_k(n; r,s) \leq T_{k-3}\left(n^{C \binom{r}{s}^2 \log \binom{r}{s}}\right).
    \end{equation*}
\end{theorem}
